\newpage
\section*{Đề bài}
\addcontentsline{toc}{section}{Đề bài}

\par
Dựa vào lý thuyết \textbf{``Chương 2 -- Analysis and interpretation of data''} đã học.
Hãy tính toán và thực hiện các yêu cầu bên dưới như sau:



\subsection*{Phần A: Dùng công cụ Excel, SPSS, Python và R tính toán.}

\textbf{1.1.} Giả sử có dữ liệu trong Bảng \ref{tab:cau1a} bên dưới và hệ số $\alpha$ là 5\%.

\begin{center}
\begin{table}[!htp]
\centering
\begin{tabular}{|c|c|c|}
\hline
\textbf{Nhóm 1} & \textbf{Nhóm 2} & \textbf{Nhóm 3} \\
\hline
6.5 & 5.2 & 6.3 \\
6.8 & 5.9 & 6.4 \\
6.3 & 5   & 6.7 \\
7.1 & 6.1 & 6.9 \\
7.8 & 7.1 & 6.1 \\
7.5 &     & 7.1 \\
7.9 &     &     \\
\hline
\end{tabular}
\caption{Số liệu}
\label{tab:cau1a}
\end{table}
\end{center}

Hãy thực hiện phân tích \textbf{One-way ANOVA} và kiểm định \textbf{Tukey’s HSD}, sau đó đưa ra kết luận từ kết quả phân tích.

\textbf{1.2. Phân tích Chi-square test}

\begin{itemize}
\item Giả sử có 100 người và có trình độ học vấn khác nhau tại TP.HCM ở file dữ liệu \texttt{Chisquare.sav} đính kèm.
\item Câu hỏi đặt ra là có sự liên quan giữa giới tính ``GIOITINH (Nam, Nữ)'' và trình độ học vấn 
``BANGCAP (Cao đẳng, Đại học, Sau đại học)'' hay không?
\item Dựa vào lý thuyết đã học phần Chi-square test tại Chương 2, hãy phân tích và đưa ra kết luận của mình.
\end{itemize}


\newpage
\subsection*{Phần B: Hãy tính toán các nội dung từng bước (bằng tay, bút mực), dựa vào các công thức được học ở Chương 2.}

\textbf{2.1.} Giả sử bạn có một ví dụ tương tự về phân phối $t$ (\textit{t-distribution}) trong phần khoảng tin cậy cho giá trị trung bình của quần thể (\textit{confidence interval for the population mean}) tại Chương 2.  
Bạn hãy đưa ra kết luận của mình và tính khoảng tin cậy với một mẫu khác gồm 25 lốp xe và mức độ tin cậy là 99\%.

\textbf{2.2.} Giả sử bạn có một ví dụ tương tự về phân phối $z$ (\textit{z-distribution}) trong phần khoảng tin cậy cho giá trị trung bình của quần thể (\textit{confidence interval for the population mean}) tại Chương 2.  
Bạn hãy đưa ra kết luận của mình và tính khoảng tin cậy với một mẫu ngẫu nhiên khác gồm 300 nhà quản lý, giá trị trung bình mẫu là 45{,}000 đô la và mức độ tin cậy là 98\%.

\textbf{2.3.} Giả sử bạn có một ví dụ tương tự về khoảng tin cậy cho tỷ lệ quần thể (\textit{confidence interval for a population proportion}) tại Chương 2.  
Bạn hãy đưa ra kết luận của mình và tính khoảng tin cậy với 600 người khảo sát khác, trong đó 450 người trả lời là có, và mức độ tin cậy là 97\%.

\textbf{2.4.} Giả sử bạn có một ví dụ tương tự vấn đề 1 (\textit{issue 1}) trong phần kiểm định giả thuyết giá trị trung bình quần thể (\textit{hypothesis test for a population mean}) tại Chương 2.  
Bạn hãy tính $|T|$ và đưa ra kết luận chấp nhận $H_0$ hoặc bác bỏ $H_0$ dựa trên các tham số mới, chẳng hạn như:
\begin{itemize}
\item Mức tăng trưởng trung bình đã biết là 1.4 inch/năm thay vì 1.35 inch/năm.
\item Mức tăng trưởng đường kính trung bình của mẫu là 1.65 inch/năm thay vì 1.6 inch/năm.
\item Mức độ tin cậy là 96\%.
\end{itemize}

\textbf{2.5.} Giả sử bạn có một ví dụ tương tự về vấn đề 3 -- trường hợp 1 (\textit{issue 3 -- case 1}) trong phần kiểm định tính bằng nhau của hai trung bình tổng thể (\textit{testing equality of two population means}) tại Chương 2.  
Bạn hãy đưa ra kết luận của mình và tính toán lại với các tham số mới, chẳng hạn như $n_1 = 214$, $n_2 = 375$ và mức ý nghĩa là 5\%.

\newpage
\subsection*{Chú ý}

\begin{itemize}
\item Kết quả phần (A) và (B) phải được tóm tắt và chụp hình đưa vào tệp tin có quy cách đặt tên:
\texttt{MSSV\_Hoten\_Lab02\_Output.docx (.pdf)}
\item Gửi đính kèm các file Excel, SPSS, R và Python của Phần A:
\begin{itemize}
\item \texttt{MSSV\_Hoten\_Lab02\_Excel.xlsx}
\item \texttt{MSSV\_Hoten\_Lab02\_SPSS.sav (.spv)}
\item \texttt{MSSV\_Hoten\_Lab02\_Python.ipynb}
\item \texttt{MSSV\_Hoten\_Lab02\_R.r}
\end{itemize}
\item Đóng gói thành 01 tập tin duy nhất:
\texttt{MSSV\_Hoten\_Lab02.rar (.zip)}
\item Deadline nộp bài: \textbf{10/03/2026}
\item Sinh viên nộp bài không đúng deadline xem như không có bài Lab02.
\item Nộp không đủ file hoặc đặt tên sai quy cách sẽ bị trừ điểm.
\item Các bài có hình thức trình bày và nội dung giống nhau sẽ nhận điểm 0.
\end{itemize}